High-capacity intrusion detectors are difficult to place on wireless sensor and edge devices, yet compression studies often stop at parameter counts or desktop accuracy. This paper presents \method{}, a controlled framework for testing curriculum-guided knowledge distillation with explainable artificial intelligence across five questions: predictive compression, training-route effects, explanation agreement, numerical deployment fidelity, and cross-dataset robustness. Under a train-only-scaler five-class \dataset{} protocol (fixed stratified split, ten seeds), a 1,189-parameter student distilled from a calibrated random forest attains $0.92034\pm0.00618$ \mf{} using 4.64~KB of FP32 parameters; a 3,397-parameter RF-KD student attains $0.93917\pm0.01225$ in 13.27~KB (KD$-$scratch $+0.0094$ and $+0.0078$, respectively). These students are 58.8$\times$ and 20.6$\times$ smaller than the full neural teacher on the same parameter-storage basis. RF-KD benefit is protocol-sensitive: feature-group-disjoint partitions remove Student A KD-over-scratch advantage. Deployment RF-KD SHAP rank agreement with the RF teacher is low and non-significant ($\rho\approx0.238/0.225$). Fixed-point RF-distilled exports reproduce every generated integer reference in four complete 56,200-record replays on ESP32-C3 and Arduino UNO R4 WiFi (agree$=1.0$); measured float-to-fixed macro-F1 drops are $\approx0.024$--$0.027$ under a disclosed 0.03 bound. Edge-IIoTset literature protocol exposes $\approx17\%$ test cross-partition exact-feature leakage under random-row splits; group-aware re-evaluation is reported separately. The evidence establishes KB-scale executable model cores while showing that predictive gain, cross-model explanation agreement, and conversion fidelity are independent outcomes.
